\section*{Getting Started}
\addcontentsline{toc}{section}{Getting Started}

\subsection*{Prerequisites and Environment Setup}

Symphony's development environment requires specific tools and configurations to support our dual-language architecture and research objectives. The following prerequisites ensure compatibility with our Rust-Python hybrid system and enable effective contribution to the research codebase.

\begin{table}[h]
\centering
\begin{tabular}{@{}lll@{}}
\toprule
\textbf{Component} & \textbf{Version} & \textbf{Purpose} \\
\midrule
Rust & 1.75+ & Core infrastructure and performance-critical components \\
Python & 3.11+ & Conductor orchestration and machine learning \\
Node.js & 18+ & Frontend development and build tools \\
Git & 2.40+ & Version control and collaboration \\
Docker & 24+ & Containerized development environment \\
\bottomrule
\end{tabular}
\caption{Development Environment Prerequisites}
\end{table}

\begin{infobox}[title=Research Environment Considerations]
Our development environment is optimized for research reproducibility and empirical validation. The specific version requirements ensure consistent behavior across different research contexts and enable reliable performance benchmarking.
\end{infobox}

\subsection*{Repository Structure and Initial Setup}

Symphony employs a monorepo architecture that facilitates integrated development across multiple research domains. The repository structure reflects our architectural principles and supports both individual component development and system-wide research activities.

\begin{expandedlist}
    \item \textbf{Clone Repository}: Access the Symphony research repository through our academic collaboration platform
    \item \textbf{Environment Configuration}: Set up development environment variables for research and debugging
    \item \textbf{Dependency Installation}: Install all required dependencies using our automated setup scripts
    \item \textbf{Verification Testing}: Run initial test suite to verify environment configuration
\end{expandedlist}

The setup process includes configuration of development tools, research data directories, and performance monitoring systems essential for empirical evaluation of architectural innovations.

\subsection*{Development Workflow Overview}

Our development workflow integrates research methodology with software engineering best practices. Contributors engage in both theoretical analysis and practical implementation, ensuring that code contributions advance our research objectives while maintaining system reliability.

\begin{description}[leftmargin=3cm,labelwidth=2.5cm]
    \item[\textbf{Research Phase}] Literature review, problem analysis, and theoretical framework development
    \item[\textbf{Design Phase}] Architectural decision-making with empirical justification
    \item[\textbf{Implementation}] Code development following research-driven specifications
    \item[\textbf{Evaluation}] Performance analysis and comparative assessment
    \item[\textbf{Documentation}] Academic-quality documentation of contributions and insights
\end{description}

This workflow ensures that every contribution advances both the practical capabilities of Symphony and the theoretical understanding of AI-first development environments.